\documentclass[t,usepdftitle=false]{beamer}

\usepackage[utf8]{inputenc}
\usetheme{Singapore}
\usepackage{xcolor}
\setbeamertemplate{footline}[frame number]

\usepackage{etex}
\usepackage{pictex}
\usepackage{tikz}
\usetikzlibrary{shapes,arrows}
\usepackage{pgfplots}

\usepackage{amsmath}

\input{../sty/macros}

% \setbeamercovered{transparent}
%\usecolortheme{crane}
\title[Simulation]{Discrete event simulation}

\author[Fabian Bastin]{Fabian Bastin \\ \url{fabian.bastin@umontreal.ca} \\ Université de Montréal -- CIRRELT -- IVADO}
\date{}

\def\bu{\boldsymbol{u}}
\def\bU{\boldsymbol{U}}
\def\cS{\mathcal{S}}

\begin{document}
\frame{\titlepage}

% ------------------------------------------------------------------------------------------------------------------------------------------------------
\begin{frame}
	\frametitle{Basic concepts}
	
	\begin{description}
		\item[System]
		A collection of entities that act and interact in order to achieve a certain logical goal.\\
		The state of a system is the collection of variables necessary to describe the system at a particular moment.
		\item[Model]
		A simplified description of a system, aimed at evaluating its performance or the effect of certain decisions.
		\item[Simulation]
		Evolving the model of a system by providing appropriate inputs, then observing and analyzing the results.
	\end{description}
	
	\mbox{}
	
	Two main system types: discrete and continuous.
	
\end{frame}

\begin{frame}
	\frametitle{Models}
	
	Physical vs mathematical model.
	
	\mbox{}
	
	Analytical model, numerical model, model with simulation.

\mbox{}

Reference textbook: 
A. M. Law, Simulation Modeling and Analysis, fifth edition, McGraw-Hill, USA, 2015.

\end{frame}

\begin{frame}
\frametitle{Event approach}

Assumption: system in which state variables can only change by a countable number of points in time.

\mbox{}

Sequence of events $e_0, e_1, e_2\ldots$, occurring at times $0 \leq t_0 \leq t_1 \leq t_2 \leq \ldots$. 

\mbox{}

${\cS_i}$: system state immediately after ${e_i}$.

\mbox{}

simulation time: current value of ${t_i}$.

\mbox{}

$(t_i,\cS_i)$ must contain enough information to continue the simulation (except the realizations of the random variables for events $e_j$, $j > i$).

\end{frame}

\begin{frame}
\frametitle{Discrete events}

\begin{center}
%%  Evolution of a general discrete-event model.
\hspace{0.5cm}
\beginpicture
\setcoordinatesystem units <7pt,7pt>
%%  72.27pt = 1in
\setplotarea x from 0 to 36, y from 0 to 10
\axis left
  label {\lines {\emph{state}\cr\vbox to 95pt{\hbox to 1.5pt{\null }} \cr}
         \hskip 9pt } /

\axis bottom
  label {\hbox to 4.0in {\hfill \emph{time}}}
  ticks length <4pt> withvalues 
     0 $t_1$ $t_2$ $t_3$ $t_4$ $t_5$ $t_6$ /
  at 0.0 6.0 7.5 12.5 18.0 25.2 30.0 / /

\axis bottom /
\put {$\rightarrow$} at 36.0 -0.1 
\put {$\uparrow$}    at 0.0 10.0 
\multiput {$\bullet$} at
  0.0 2.0
  6.0 1.6
  7.5 2.7
 12.5 6.1
 18.0 5.1
 25.2 2.6
 30.0 4.0 /
\putrule from  0.0 2.0 to  6.0 2.0
\putrule from  6.0 1.6 to  7.5 1.6
\putrule from  7.5 2.7 to 12.5 2.7
\putrule from 12.5 6.1 to 18.0 6.1
\putrule from 18.0 5.1 to 25.2 5.1
\putrule from 25.2 2.6 to 30.0 2.6
\putrule from 30.0 4.0 to 32.0 4.0 
\endpicture
\end{center}

\end{frame}

\begin{frame}
\frametitle{General framework}

\begin{itemize}
\item
Main program.
\item
Simulation clock.
\item
Event list.
\item
Random numbers generation functions.
\item
Statistical collectors.
\item
Report generation.
\end{itemize}

Event: action, schedule, cancel.

\end{frame}

\begin{frame}

\tikzstyle{decision} = [diamond, draw, fill=blue!20,
    text width=2.5cm, text badly centered, node distance=1cm, inner sep=0pt]
\tikzstyle{block} = [rectangle, draw, fill=blue!20,
    text width=2.5cm, text centered, rounded corners, minimum height=4em]
\tikzstyle{line} = [draw, very thick, color=black!50, -latex']
\tikzstyle{cloud} = [draw, ellipse,fill=yellow!20, node distance=1cm,
    minimum height=2em]

\begin{footnotesize}
\begin{center}
\begin{tikzpicture}[scale=0.09, node distance = 0.8cm, auto]
    % Place nodes
    \node [block] (init) {Initialization};
    \node [text width=2.2cm, cloud, left of=init, node distance=3.7cm] (clock) {Set simulation clock to 0};
    \node [text width=2.5cm, cloud, right of=init, node distance=3.7cm] (system) {Plan first event(s)};
    \node [text width=2.1cm, cloud, below of=init, node distance=2.0cm] (premier) {Pull the first event};
    \node [text width=2.1cm, block, below of=premier, node distance=2.0cm] (event) {Handle current event};
    \node [text width=2.1cm, block, below of=event, node distance=2.0cm] (schedule) {Plan resulting events};
    \node [block, left of=schedule, node distance=4.0cm] (update) {Pull next event and advance clock};
    \node [text width=2.1cm, decision, below of=schedule, node distance=4.0cm] (decide) {Liste d'événements vide ou événement de fin?};
    \node [block, below of=decide, node distance=4.5cm] (stop) {Stop};
    % Draw edges
    \path [line] (init) -- (premier);
    \path [line] (premier) -- (event);
    \path [line] (event) -- (schedule);
    \path [line] (schedule) -- (decide);
    \path [line] (decide) -| node [near start, color=black] {non} (update);
    \path [line] (update) |- (event);
    \path [line] (decide) -- node [, color=black] {oui}(stop);
    \path [line,dashed] (clock) -- (init);
    \path [line,dashed] (system) -- (init);
\end{tikzpicture}
\end{center}
\end{footnotesize}

\end{frame}

\begin{frame}

\tikzstyle{decision} = [diamond, draw, fill=blue!20,
    text width=2.5cm, text badly centered, node distance=2.5cm, inner sep=0pt]
\tikzstyle{block} = [rectangle, draw, fill=blue!20,
    text width=2.5cm, text centered, rounded corners, minimum height=4em]
\tikzstyle{line} = [draw, very thick, color=black!50, -latex']
\tikzstyle{cloud} = [draw, ellipse,fill=yellow!20, node distance=2.5cm,
    minimum height=2em]

\begin{footnotesize}
\begin{center}
\begin{tikzpicture}[scale=0.09, node distance = 1cm, auto]
    % Place nodes
    \node [text width=3cm, block, below of=premier, node distance=2cm] (event) {Handle current event};
    \node [text width=3cm, block, below of=event, node distance=2cm] (schedule) {Plan resulting events};
    \node [block, left of=schedule, node distance=4.5cm] (update) {Pull next event and advance clock};
    \node [text width=1.9cm, decision, below of=schedule, node distance=2.5cm] (decide) {End event or empty event list?};
    \node [block, below of=decide, node distance=2.5cm] (stop) {Stop and print report};
    % Draw edges
    \path [line] (event) -- (schedule);
    \path [line] (schedule) -- (decide);
    \path [line] (decide) -| node [near start, color=black] {no} (update);
    \path [line] (update) |- (event);
    \path [line] (decide) -- node [, color=black] {yes}(stop);
\end{tikzpicture}
\end{center}
\end{footnotesize}

\end{frame}

\begin{frame}
\frametitle{Queueing systems}

\includegraphics[width=\textwidth]{queue-labels.png}

\end{frame}

\begin{frame}
\frametitle{Kendall notation}

$$
A/S/c/K/N/SD
$$

\mbox{}

\begin{description}
	\item[$A$]
Arrival process
	\item[$S$]
Service time distribution
	\item[$c$]
Number of servers
	\item[$K$]
Waiting list maximum length
	\item[$N$]
Population size
	\item[$SD$]
Service discipline
\end{description}
 
\end{frame}

\begin{frame}
\frametitle{Notation}

\begin{description}
	\item[$G$]
	General
	\item[$I$]
	Independent
	\item[$M$]
	Markovian or memoryless
\end{description}

\mbox{}

Defaults:
\begin{itemize}
	\item
	Infinite buffer capacity
	\item
	Infinite population size
	\item
	FCFS (First-Come-First-Served) service discipline
\end{itemize}

\mbox{}

Memoryless property:
$$
P[X > x +a \,|\, X > a] = P[X > x]
$$
Continuous: exponential distribution\\
Discrete: geometric distribution 

\end{frame}

\begin{frame}
\frametitle{Example: GI/G/1 queue}

$$
G/G/1 = G/G/1/\infty/\infty/FCFS
$$

We also write $GI/G/1$ to stress that the service times are independent.

\mbox{}

The clients arrive one by one, one being served at a time, under a FCFS priority, with
\begin{itemize}
\item
${S_i}$: service time of client $i$, with cdf $G$;
\item
${A_i}$: inter-arrival time between clients $i$ and $i+1$, with cdf $F$.
\end{itemize}
We assume $S_i$ and $A_i$ mutually independent.

\end{frame}

\begin{frame}
	\frametitle{Example: GI/G/1 queue}
	
	The first client arrives at time ${A_0}$ in an empty system.
	
	\mbox{}
	
	We wish to simulate this system for a duration ${T}$ and calculate the average waiting time per client and the average queue length. The types of events are
	\begin{enumerate}
		\item Arrival;
		\item Departure;
		\item End of simulation.
	\end{enumerate}
	The random variables (independent) to generate are $A_1, A_2, A_3, \dots$ and $S_1, S_2, S_3, \dots$.
	
\end{frame}

\begin{frame}
	\frametitle{Example: GI/G/1 queue}
	
	\begin{itemize}
		\item ${W_i}$, the waiting time of client $i$;
		\item ${Q(t)}$, the length of the queue at time $t$;
		\item ${N_c(t)}$, the number of clients who started their service during the interval $[0,t]$.
	\end{itemize}
	Suppose we want to calculate, for the interval $[0,T]$, the \emph{average waiting time} per client,
	$${\overline{W}_{N_c(T)}} = {1\over N_c(T)} \sum_{i=1}^{N_c(T)} W_i;$$
	and the \emph{average queue length}:
	$${\overline{Q}_T} = {1\over T} \int_0^T Q(t) dt.$$
	
\end{frame}

\begin{frame}
	\frametitle{Example: GI/G/1 queue}
	
	This last integral is easy to calculate through simulation.
	
	\mbox{}
	
	For each \emph{client}, we create an \emph{object} containing the arrival time and service duration.
	
	\mbox{}
	
	State variables: list of clients waiting and the list of clients in service.
	
	\mbox{}
	
	+ simulation clock, statistical counters (if needed), \emph{list} of future scheduled events, in chronological order, and a procedure for each type of event.
	
\end{frame}

\begin{frame}
	\frametitle{Queueing System Simulation}
	
	\begin{small}
		\begin{description}
			\item[Arrival]
			Generate $A$ according to $F$ and schedule another arrival 
			in $A$ time units;\\
			Create the new client and note their arrival time;\\
			Generate their service duration $S$ according to $G$;\\
			{\bf If} (server is busy) {\bf then}\\
			\quad Insert this client into the waiting list;\\
			{\bf else}\\
			\quad Insert the client into the service list;\\
			\quad Schedule their departure in $S$ time units; \\
			Update the desired statistics.
			\item[Departure]
			Remove the client from the service list;\\
			{\bf If} the queue is not empty {\bf then}\\
			\quad Remove the first client from the waiting list;\\
			\quad Insert them into the service list;\\
			\quad Schedule their departure in 
			$S$ time units;\\
			Update the desired statistics.
			\item[End of Simulation]
			Print a report and terminate the program.
		\end{description}
	\end{small}
\end{frame}

\begin{frame}
	\frametitle{Simulation}
	
	Start the simulation:
	\begin{itemize}
		\item Initialize variables and counters,
		\item Schedule the ``End of Simulation'' at time $T$,
		\item Generate $A$ according to $F$ and schedule the first ``Arrival'' in $A$ time units,
		\item Launch the simulation.
	\end{itemize}
	
	\mbox{}
	
	The execution of the simulation simply consists of repeating the following loop:\\
	{\bf Repeat:} execute the next event
	in the event list\\
	{\bf until:} the list of scheduled events is empty,\\
	\quad    or an event stops the simulation.
	
\end{frame}

\begin{frame}
	\frametitle{Simulation without Event List}
	
	It is not always necessary to resort to the event-driven approach.
	
	\mbox{}
	
	Example: Lindley's recurrence:
	\[
	W_1 = 0,\quad W_{i+1} = \max(0,\; W_i + S_i - A_{i+1}).
	\]
	
	We can easily simulate for a fixed number of clients (instead of a fixed time horizon $T$).
	
	\mbox{}
	
	The program formulation based on Lindley's recurrence reduces to solving an integration problem over $[0,1)^t$.
	
\end{frame}

\begin{frame}
	\frametitle{Lindley's Recurrence}
	
	Suppose we want to estimate \(E[\overline{W}_{100}]\) by \({\overline{W}_{100}}\):
	\[
	{\overline{W}_{100}} = \frac{1}{100} \sum_{i=1}^{100} W_i.
	\]
	
	\mbox{}
	
	We need \(A_1, S_1, A_2, S_2, \dots, A_{99}, S_{99}\). If we set \({A_i} = F^{-1}(U_{2i-1})\) and \({S_i} = G^{-1}(U_{2i})\), where \(U_j\), \(j = 1,\ldots,99\), are uniform random variables over \((0,1)\), then \(\overline{W}_{100} = f(U_1,\dots,U_{198})\) (with \(f\) being an unknown function in this case).
	
	\mbox{}
	
	We then have \(t=198\), i.e., 
	\[
	E[\overline{W}_{100}] = \int_{[0,1)^{198}} f(\bu)d\bu.
	\]
	
\end{frame}

\begin{frame}
	\frametitle{Random Number of Clients}
	
	On the other hand, if we want to simulate \(\overline{W}_{N_c(T)}\), the number of clients 
	\(N_c(T)\) is random and unbounded. 
	We should therefore choose \(t=\infty\) as the dimension.
	
\end{frame}

\begin{frame}
	\frametitle{Process Approach}
	
	A process is a temporally ordered sequence of interrelated events, separated by certain time intervals,
	that describes the entire experience of an ``entity'' as it evolves through a ``system''.
	A simulation system or model may involve different types of processes.
	
	\mbox{}
	
	A routine will be associated with each process in the model, describing its
	entire history through the system.
	
	\mbox{}
	
	In contrast to the event-driven approach, a process routine
	explicitly contains the passage of simulated time and generally has
	multiple entry points.
	
\end{frame}

\begin{frame}
	\frametitle{Processes vs Events}
	
	A simulation using the process approach also evolves over time by executing events in the order of their occurrence.
	
	\mbox{}
	
	Internally, the event-driven and process approaches are very similar.
	
	\mbox{}
	
	For the user, reasoning may differ, and it can sometimes (but not always) feel more natural to use the process approach.
	
	%\mbox{}
	
	%Moreover, the process approach is often slower in terms of execution time, as the simulator ultimately has to handle events and manage potentially concurrent processes.
	
\end{frame}

\begin{frame}
	\frametitle{Improving Efficiency}
	
	The following material is adapted from Pierre L'Ecuyer's class ``Simulation: stochastic aspects''.
	
	\mbox{}
	
	The efficiency of an estimator \(X\) is defined as follows:
	$$
	eff[X] = \frac{1}{MSE[X] C(X)},
	$$
	where MSE stands for mean square error and \(C(X)\) represents the cost to compute \(X\).
	
	\mbox{}
	
	Is it possible, given \(X\), to construct a new estimator \(Y\) that is more efficient than \(X\)?
	
	\mbox{}
	
	We will introduce the main concepts of improving efficiency using an introductory example about call centers.
	
\end{frame}

\begin{frame}
	\frametitle{Introductory Example}
	
	Let \(B\) be the \emph{load factor} for the day, and suppose that \(P[B = b_t] = q_t\), where
	\begin{center}
		\begin{tabular}{c|cccc}
			${t}$   &   1  &  2   &   3  &  4  \\  
			\hline
			${b_t}$ &  0.8 &  1.0 &  1.2 &  1.4 \\
			${q_t}$ & 0.25 & 0.55 & 0.15 & 0.05 \\
		\end{tabular}
	\end{center}
	
	\mbox{}
	
	It is easy to verify that \(E[B] = 1\). %% and \(\Var[B] = 0.024\).
	
	\mbox{}
	
	The \emph{arrivals} of calls follow a Poisson process with rate
	\(B{\lambda_j}\) during hour \(j\).
	
	\mbox{}
	
	Let \(G_i(s)\) be the number of calls that started service
	after less than \(s\) seconds of waiting on day \(i\).
	
	\mbox{}
	
	We want to estimate \(\mu = E[G_i(s)]\), say for \(s=20\).
	
\end{frame}

\begin{frame}
	\frametitle{Example}
	
	We also assume that the service durations of the calls follow the distribution 
	\(\Gamma(\alpha,\gamma)\), with a mean of \(\alpha\gamma\).
	
	\mbox{}
	
	In our example, we have \(\alpha = 1\) and \(\gamma = {\gamma_1} = 100\).
	
	\mbox{}
	
	Number of agents \(n_j\) and arrival rate \(\lambda_j\) (per hour) for 13 hourly periods in the call center:
	
	\begin{footnotesize}
		\hspace*{-1cm}
		\begin{tabular}{l|rrrrrrrrrrrrr}
			\hline
			$j$   & 0 & 1 & 2 & 3 & 4 & 5 & 6 & 7 & 8 & 9 & 10 & 11 & 12 \\
			$n_j$ & 4 & 6 & 8 & 8 & 8 & 7 & 8 & 8 & 6 &  6 &  4 &  4 &  4 \\
			$\lambda_j$ & 100 & 150 & 150 & 180 & 200 & 150 & 150 & 150 & 
			120 & 100 & 80 & 70 & 60 \\
			\hline
		\end{tabular}
	\end{footnotesize}
	
	\mbox{}
	
	We simulate \(n\) independent days.
	
\end{frame}

\begin{frame}
	\frametitle{Example}
	
	Soit ${X_i} = G_i(s)$ pour le jour $i$, et 
	\[
	{\bar X_n} = {1\over n} \sum_{i=1}^n X_i.
	\]
	On a $E[\bar X_n] = \mu$ et $Var[\bar X_n] = Var[X_i]/n$.
	
	\mbox{}
	
	Une expérience avec ${n = 1000}$ donne $\bar X_n = 1518.3$ et $S_n^2
	= {21615}$.
	La variance estimée de $\bar X_n$ est alors $\widehat{Var}[\bar X_n] = 21.6$.
	
\end{frame}

\begin{frame}
	\frametitle{Introductory Example}
	
	We want to construct a confidence interval for \(\sigma^2 = n Var[\bar X_n]\), under the assumption that \((n-1)S_n^2/\sigma^2\) approximately follows a \(\chi^2\) distribution with \(n-1\) degrees of freedom. This allows us to construct the \(90\%\) interval: \([0.930S_n^2, 1.077S_n^2]\).
	
	\mbox{}
	
	In other words, the relative error of this estimator is less than \(8\%\) with a ``confidence'' of approximately \(90\%\).
	
	\mbox{}
	
	Let's see how to \emph{improve} this estimator \(\bar X_n\) by reducing its variance. For each proposed method, we will provide numerical results for \(n = 1000\).
	
\end{frame}

\begin{frame}
	\frametitle{Indirect Estimation}
	
	Let \(A_i\) be the total number of arrivals on day \(i\); let \(D_i = A_i - X_i\).
	
	\mbox{}
	
	We know that \(a = E[A_i] = \sum_{j=1}^m \lambda_j = 1660\).
	
	\mbox{}
	
	Thus, we can write \(\mu = E[X_i] = E[A_i - D_i] = a - E[D_i]\), which 
	can be estimated by
	$$ {\bar X_{\rm{i},n}} = E[A_i] - \bar D_n = a - {1\over n} \sum_{i=1}^n D_i.$$
	
	\mbox{}
	
	This estimator has less variance than \(\bar X_n\) if and only if \(Var[D_i] < Var[X_i]\). 
	Estimated variance: \(\widehat{Var}[X_{\rm{i},i}] = {18389}\).
	
\end{frame}

\begin{frame}
	\frametitle{Control Variate (CV)}
	
	Idea: leverage auxiliary information.
	
	\mbox{}
	
	For example, if \(A_i\) is higher than usual \((A_i > E[A_i] = 1660)\),
	we expect that on that day, \(X_i\) and \(D_i\) 
	\emph{overestimate} \(E[X_i]\) and \(E[D_i]\).
	
	\mbox{}
	
	We could make a ``\emph{correction}'' to these estimators:
	replace \(X_i\) with
	$$ X_{\rm{c},i} = X_i - {\beta} (A_i - 1660) $$
	where \(\beta\) is an appropriate constant. Then
	\[
	{\bar X_{\rm{c},n}} = \bar X_n - \beta (\bar A_n - 1660).
	\]
	
\end{frame}

\begin{frame}
	\frametitle{Control Variate (CV)}
	
	We have \(E[\bar X_{\rm{c},n}] = E[X_i]\) and 
	$$ Var[\bar X_{\rm{c},n}] 
	= \frac{Var[X_i] + \beta^2Var[A_i] - 2 \beta Cov[A_i,X_i]}{n}. $$ 
	This variance is a quadratic function of \(\beta\), which we minimize
	by taking
	\[
	\beta = {\beta^*} = \frac{Cov[A_i, X_i]}{Var[A_i]}.
	\]
	
	\mbox{}
	
	The \emph{minimum variance} is
	$$ Var[\bar X_{\rm{c},n}] = \frac{Var[X_i] - (\beta^*)^2 Var[A_i]}{n}
	= Var[\bar X_n] (1- \rho^2[A_i,X_i]) $$
	where \({\rho[A_i, X_i]}\) is the correlation coefficient between 
	\(A_i\) and \(X_i\).
	
\end{frame}

\begin{frame}
	\frametitle{Control Variate (CV)}
	
	We do not know \(Cov[A_i, X_i]\), but:
	\begin{enumerate}[(a)]
		\item We can estimate it through \emph{pilot experiments}.
		\item We can estimate it using the \emph{same \(n\) simulations} as \(\bar X_n\).
	\end{enumerate}
	
	\mbox{}
	
	With (b), we obtain the estimator (slightly biased):
	\[
	\bar X_{\rm{c},e,n} = \bar X_n - {1\over n-1} \left[\sum_{i=1}^n
	(A_i-\bar A_n)(X_i-\bar X_n)\right] {\bar A_n - a\over Var[A_i]}.
	\]
	
	\mbox{}
	
	Conditional on \(B_i\), \(A_i \sim \text{Poisson}(1660 B_i)\). Thus, we have
	\begin{align*}
		{Var[A_i]} &= Var[E[A_i|B_i]] + E[Var[A_i|B_i]] \\
		&= Var[1660 B_i] + E[1660 B_i] \\
		&= 1660^2 Var[B_i] + 1660 E[B_i] \\
		&= 67794.4.
	\end{align*}
	The empirical variance obtained here is: {3310}.
	
\end{frame}

\begin{frame}
	\frametitle{Control Variable (CV)}
	
	By taking \(\beta=1\), we recover the indirect estimator:
	\[
	\bar X_{\rm{i},n} = \bar X_n - (\bar A_n - 1660) = 1660 - \bar D_n.
	\]
	If we combine \emph{CV + indirect}, we obtain:
	\begin{eqnarray*}
		{\bar X_{\rm{i},\rm{c},n}} &=& a - \bar D_n - \beta_2 (\bar A_n - a) \\
		&=& \bar A_n - \bar D_n -(1+\beta_2)(\bar A_n - a) \\
		&=& \bar X_n -(1+\beta_2)(\bar A_n - a),
	\end{eqnarray*}
	i.e., \(\bar X_{\rm{i},\rm{c},n}\) is equivalent to \(\bar X_{\rm{c},n}\) 
	with \(\beta = 1+\beta_2\).
	As a result, in the presence of the control variable, the indirect estimation provides no additional benefit.
	
	\mbox{}
	
	We could also consider other control variables, such as \(B_i\), the average service durations, etc.
	
\end{frame}

\begin{frame}
	\frametitle{Stratification}
	
	The load factor \(B_i\) is a significant source of variability in this case. Let's try to control it.
	
	\mbox{}
	
	By setting \(\mu_t = E[X_i\mid B_i=b_t]\), we have
	\begin{eqnarray*}
		\mu &=& E[X_i]  = \sum_{t=1}^4 P[B_i = b_t] \cdot E[X_i\mid B_i=b_t] \\
		&=&  0.25\, E[X_i\mid B_i=0.8] + 0.55\, E[X_i\mid B_i=1.0]  \\
		&& + 0.15\, E[X_i\mid B_i=1.2] + 0.05\, E[X_i\mid B_i=1.4]  \\
		&=&  0.25\,\mu_1 + 0.55\, \mu_2 + 0.15\, \mu_3 + 0.05\, \mu_4.
	\end{eqnarray*}
	
	\mbox{}
	
	\emph{Idea}: estimate \(\mu_t\) \emph{separately} 
	for each \(t\).
	
\end{frame}

\begin{frame}
	\frametitle{Stratification}
	
	Suppose there are \(N_t\) days where \(B_i = b_t\) and let
	\(X_{t,1},\dots, X_{t,N_t}\) be the values of \(X_i\) for those days. 
	
	\mbox{}
	
	We can estimate \(\mu_t = E[X_i\mid B_i=b_t]\) by
	$$ {\hat\mu_t} = {1\over N_t} \sum_{i=1}^{N_t} X_{t,i} $$
	and \(\mu\) by
	\[
	{\bar X_{\rm{s},n}} = \sum_{t=1}^4 q_t \hat\mu_t 
	= 0.25 \hat\mu_1 + 0.55 \hat\mu_2 + 0.15 \hat\mu_3 + 0.05 \hat\mu_4.
	\]
	
\end{frame}

\begin{frame}
	\frametitle{Stratification}
	
	We have 
	\begin{align*}
		&  Var[\bar X_{\rm{s},n} \mid N_1, N_2, N_3, N_4] \\
		&= \sum_{t=1}^4 q_t^2 Var[\hat\mu_t | N_t]  
		= \sum_{t=1}^4 q_t^2 \frac{\sigma^2_t}{N_t} \\
		&= 0.25^2 \frac{\sigma^2_1}{N_1} + 0.55^2 \frac{\sigma^2_2}{N_2} + 
		0.15^2 \frac{\sigma^2_3}{N_3} + 0.05^2 \frac{\sigma^2_4}{N_4}.
	\end{align*}
	where \(\sigma^2_t = Var[X_i\mid B_i=b_t]\).
	
	\mbox{}
	
	The variance is reduced if the \(\sigma_t^2\) are less than \(Var[X_i]\).
	
	\mbox{}
	
	If the \(B_i\) are normally generated: \emph{post-stratification}.
	
	\mbox{}
	
	To estimate \(\mu\) through stratification, we can also 
	\emph{fix the \(N_t = n_t\) in advance}, meaning we choose 
	in advance how many days will have \(B_i = b_t\) for each value of 
	\(t\).
	
\end{frame}

\begin{frame}
	\frametitle{Stratification}
	
	\begin{itemize}
		\item
		\emph{Proportional allocation}: set \(n_t = n q_t\).\\
		With \(n=1000\), this gives \(n_1 = 250\), \(n_2 = 550\), \(n_3 = 150\), \(n_4 = 50\).
		\item
		\emph{Optimal allocation}: 
		choose the \(n_t\) to minimize \(Var[\bar X_{\rm{s},n}]\) under the constraint
		\(n_1+n_2+n_3+n_4=n\). We obtain:
		\[
		\frac{{n_t}}{n} 
		= \frac{\sigma_t P[B_i=t]}{\sum_{k=1}^4 \sigma_k P[B_i=k]}
		= \frac{\sigma_k q_k}{\sum_{k=1}^4 \sigma_k q_k}.
		\]
		We do not know these \(\sigma_k\), but we can estimate them through pilot trials.
		With \(n_0=800\) pilot trials, 200 for each value of \(t\), we obtain, for example, \((n_1, n_2, n_3, n_4) = (219, 512, 182, 87)\) (after rounding). 
	\end{itemize}
	
\end{frame}

\begin{frame}
	\frametitle{Combined Strategies}
	
	\emph{Stratification combined with CV:}\\
	$$ 
	{\hat\mu_t} = {1\over n_t} \sum_{i=1}^{n_t} X_{\rm{c},t,i} 
	= {1\over n_t} \sum_{i=1}^{n_t} 
	[X_{t,i} - \beta_t (A_{t,i} - a\,b_t)].
	$$
	
	\mbox{}
	
	We minimize \(\sigma_t^2 = Var[X_{\rm{c},t,i}]\) by taking 
	\[
	\beta_t = {\beta_t^*} = \frac{Cov[A_{t,i},\, X_{t,i}]}{Var[A_{t,i}]}
	= \frac{Cov[A_i, X_i | B_i=b_t]}{(a b_t)}.
	\]
	
	\mbox{}
	
	The addition of a CV alters the \(\sigma_t^2\): the optimal allocation 
	is no longer the same.
	
	\mbox{}
	
	With \(n_0=800\) pilot trials, we obtain
	\((\beta_1,\beta_2, \beta_3, \beta_4) = (1.020, 0.648, 0.224, -0.202)\) 
	and \((n_1, n_2, n_3, n_4) = (131, 503, 247, 119)\) 
	as estimates of the optimal values.
	
\end{frame}

\begin{frame}
	\frametitle{Combined Strategies}
	
	By repeating the experiment with \(n = 100000\), we can find the following estimates for the variance (\(\pm 1\%\)):\\
	$Var[X_{n}] = 21998$; \
	$Var[X_{\rm{i},n}] = 17996$; \
	$Var[X_{\rm{c},n}] = 3043$; \
	$Var[X_{\rm{s}o,\rm{c},n}] = 885$.
	
\end{frame}

\begin{frame}
	\frametitle{Numerical Results for \(n = 1000\)}
	
	\begin{small}
		\begin{tabular}{llrrr}
			\hline
			Method                & Estimator
			& Mean   & $S_n^2 (\pm 9\%)$ & Ratio \\
			\hline
			&&&&\\
			Crude estimator             & $\bar X_n$
			& 1518.2 & 21615  &  1.000 \\
			Indirect                    & $\bar X_{\rm{i},n}$
			& 1502.5 & 18389  &  0.851 \\
			CV $A_i$, with pilot runs   & $\bar X_{\rm{c},n}$
			& 1510.1 &  3305  &  0.153 \\
			CV $A_i$, no pilot runs     & $\bar X_{\rm{c}e,n}$
			& 1510.2 &  3310  &  0.153 \\
			Indirect + CV, no pilot runs & $\bar X_{\rm{i},\rm{c},n}$
			& 1510.1 &  3309  &  0.153 \\
			Stratification (proportional)   & $\bar X_{\rm{s}p,n}$
			& 1509.5 &  1778  &  0.082 \\
			Stratification (optimal)    & $\bar X_{\rm{s}o,n}$
			& 1509.4 &  1568  &  0.073 \\
			&&&& \\
			Strat.\ (proportional) + CV & $\bar X_{\rm{s}p,\rm{c},n}$
			& 1509.2 &  1140  &  0.053 \\
			Strat.\ (optimal) + CV & $\bar X_{\rm{s}o,\rm{c},n}$
			& 1508.3 &   900  &  0.042 \\
			\hline
		\end{tabular}
	\end{small}
	
	%comparison of two systems
\end{frame}

\begin{frame}
	\frametitle{Comparison of Two Similar Systems}
	
	Now suppose it is possible to slightly decrease the parameter \(\gamma\) of the distribution governing the service durations, by taking
	\[
	\gamma = {\gamma_2} = \gamma_1(1-\delta),
	\]
	for \({\delta} > 0\) being very small.
	
	\mbox{}
	
	If we use the same random numbers, this is equivalent to multiplying the service durations by \((1-\delta)\).
	
	\mbox{}
	
	We want to estimate \(\mu(\gamma_2) - \mu(\gamma_1) = \EE_\gamma[X_2] - \EE_\gamma[X_1]\),
	in order to measure the effect of slightly increasing the speed of the servers.
	
\end{frame}

\begin{frame}
	\frametitle{Comparison of Two Similar Systems}
	
	We simulate \(n\) days for each value of \(\gamma\).
	\begin{itemize}
		\item
		${X_{1,i}}=$ value of \(G_i(s)\) with \(\gamma_1\);
		\item
		${X_{2,i}}=$ value of \(G_i(s)\) with \(\gamma_2\);
		\item
		${\Delta_i} = X_{2,i}-X_{1,i}$,
		\[
		{\bar\Delta_n} = {1\over n} \sum_{i=1}^n \Delta_i
		\]
	\end{itemize}
	
	\mbox{}
	
	We can simulate \(X_{1,i}\) and \(X_{2,i}\)
	\begin{enumerate}[(i)]
		\item with independent random variables (\emph{IRV}),
		\item with common random variables (\emph{CRV}).
	\end{enumerate}
	
\end{frame}

\begin{frame}
	\frametitle{Comparison of Two Similar Systems}
	
	As
	\[
	Var[\Delta_i] = Var[X_{1,i}] + Var[X_{2,i}] - 2Cov[X_{1,i},\,X_{2,i}],
	\]
	the goal of common random variables is to make the covariance positive.
	
	\mbox{}
	
	How to implement CRNs?
	
	\mbox{}
	
	Use different random sequences to generate:
	\begin{enumerate}
		\item The load factor \(B_i\);
		\item The inter-arrival times;
		\item The call durations;
		\item The patience durations.
	\end{enumerate}
	Everything is generated by inversion: a uniform distribution is used for each random variable.
	
\end{frame}

\begin{frame}
	\frametitle{Synchronization}
	
	When service durations change, waiting times also change,
	and abandonment decisions may consequently change.
	
	\mbox{}
	
	If we generate service durations only for clients who do not abandon, we may lose synchronization:
	we might have to generate fewer service durations in one system than in the other.
	
	\mbox{}
	
	Service durations can be generated:
	%
	\begin{itemize}
		\item[(a)] 
		for all calls, including those that abandon,
		\item[(b)] 
		only for served calls.
	\end{itemize}
	
\end{frame}

\begin{frame}
	\frametitle{Synchronization}
	
	Similarly, the patience duration does not need to be generated for 
	clients who do not wait. It can be generated:
	\begin{itemize}
		\item[(c)] 
		for all calls,
		\item[(d)] 
		only if necessary.
	\end{itemize}
	
\end{frame}

\begin{frame}
	\frametitle{Synchronization: results (with ${n=10^4}$)}
	
	\begin{footnotesize}
		\begin{center}
			\begin{tabular}{|l|rr|rr|rr|}
				\hline
				% &&&&&&\\
				\ Method
				& \multicolumn{2}{|c}{$\delta = 0.1$}
				& \multicolumn{2}{|c}{$\delta = 0.01$}
				& \multicolumn{2}{|c|}{$\delta = 0.001$} \\
				& $\bar \Delta_n$ & $\widehat{Var}[\Delta_i]$
				& $\bar \Delta_n$ & $\widehat{Var}[\Delta_i]$
				& $\bar \Delta_n$ & $\widehat{Var}[\Delta_i]$\\
				\hline
				IRN (a + c)                      &  55.2 & 56913 &  4.98 & 45164 &  0.66 & 44046 \\
				IRN (a + d)                      &  52.2 & 54696 &  7.22 & 45192 & -1.82 & 45022 \\
				IRN (b + c)                      &  50.3 & 56919 &  9.98 & 44241 &  1.50 & 45383 \\
				IRN (b + d)                      &  53.7 & 55222 &  5.82 & 44659 &  1.36 & 44493 \\
				CRN, no sync. (b + d)            &  56.0 &  3187 &  5.90 &  1204 &  0.19 &   726 \\
				CRN (a + c)                      &  56.4 &  2154 &  6.29 &    37 &  0.62 &   1.8 \\
				CRN (a + d)                      &  55.9 &  2161 &  6.08 &   158 &  0.74 &  53.8 \\
				CRN (b + c)                      &  55.8 &  2333 &  6.25 &   104 &  0.63 &   7.9 \\
				CRN (b + d)                      &  55.5 &  2323 &  6.44 &   143 &  0.59 &  35.3 \\
				% CRN (a + c), str.\ (opt.) & &&&&& \\ 
				% CRN (a + d), str.\ (opt.) & &&&&& \\
				% CRN (b + c), str.\ (opt.) & &&&&& \\
				% CRN (b + d), str.\ (opt.) & &&&&& \\
				\hline
			\end{tabular}
		\end{center}
	\end{footnotesize}
	
\end{frame}

\begin{frame}
	\frametitle{Induction of Correlation}
	
	Sufficient conditions for $Cov[X,Y]$ 
	to be positive or negative?
	
	\mbox{}
	
	How to maximize or minimize correlation for given marginal distributions?
	
	\mbox{}
	
	\begin{theorem}[Fréchet bounds]
		Among pairs of random variables $(X,Y)$ with marginal 
		distribution functions $F$ and $G$, the pair $(X,Y) = (F^{-1}(U), G^{-1}(U))$ 
		where $U \sim U(0,1)$, \emph{maximizes} $\rho[X, Y]$, \
		and the pair $(X,Y) = (F^{-1}(U), G^{-1}(1-U))$
		\emph{minimizes} $\rho[X, Y]$.
	\end{theorem}
	
\end{frame}

\begin{frame}
	\frametitle{Correlation Induction}
	
	For example, if the distribution function of a service duration is 
	${F}$ in the first system and ${G}$ in the second,
	and if we generate the service durations by
	$X = F^{-1}(U)$ and $Y = G^{-1}(U)$, then $Cov[X, Y] \ge 0$.
	
	\begin{theorem}
		Let \ ${X} = {f}(U)$ where ${U} \sim U(0,1)$ and \
		${Y} = {g}(V)$ where ${V} \sim U(0,1)$. Then,
		\begin{itemize}
			\item
			if $f$ and $g$ are monotonic in the same direction and $V=U$,
			then $Cov[X, Y] \ge 0$;
			\item
			if $f$ and $g$ are monotonic in the same direction and $V=1-U$,
			then $Cov[X, Y] \le 0$.
		\end{itemize}
	\end{theorem}
	
\end{frame}

\begin{frame}
	\frametitle{Common Random Numbers (CRN)}
	
	We use ${\Delta} = {X_2}-{X_1}$ to estimate ${\mu_2}-{\mu_1} = E[X_2] - E[X_1]$.
	Then
	\[
	Var[\Delta] = Var[X_1] + Var[X_2] - 2 \, Cov[X_1,X_2].
	\]
	\emph{Goal}: to induce a positive correlation between $X_1$ and $X_2$
	without changing their individual distributions.
	
	\mbox{}
	
	Technique: use the same random numbers to simulate
	both systems, trying to maintain \emph{synchronization}
	as best as possible.
	
	\mbox{}
	
	If $X_k = {f_k}(\bU_k) = f_k({U_{k,1}},{U_{k,2}},\dots)$ 
	for $k=1,2$, using RNGs everywhere means taking $\bU_1=\bU_2$.
	
\end{frame}

\begin{frame}
	\frametitle{Common Random Numbers (CRN)}
	
	\begin{theorem}
		If $f_1$ and $f_2$ are monotonic in the same direction with respect to
		all their arguments, then  $Cov[X_1, X_2] \ge 0$.
	\end{theorem}
	
	\mbox{}
	
	To \emph{maximize} the correlation, one should generate
	$X_1$ and $X_2$ directly through inversion!
	
\end{frame}
\begin{frame}
	\frametitle{Example: Lindley Process}
	
	${W_{i+1}} = \max(0,\, W_i + S_i - A_i)$, where $S_i - A_i$ is
	independent of $W_i$.
	
	\mbox{}
	
	If ${X_1}$ and ${X_2}$ are non-decreasing functions of the $W_i$
	for two Lindley processes simulated with common random numbers,
	then $Cov[X_1, X_2] \ge 0$.
	
	\mbox{}
	
	Sometimes it is very difficult to check the conditions of the
	theorem.
	
	\mbox{}
	
	For example, in a bank or a call center, monotonicity
	is hard to verify due to abandonment possibilities.
	Moreover, the conditions of the theorem are \emph{sufficient}, but 
	\emph{not necessary}.
	
\end{frame}

\begin{frame}
	\frametitle{Common Random Numbers (CRN)}
	
	To \emph{test} if it is effective in practice:
	conduct a \emph{pilot experiment with common random numbers} and estimate 
	$Cov(X_1, X_2)$.
	
	\mbox{}
	
	It is not necessary to perform simulations without common random numbers:
	to estimate the variance one would have without common random numbers,
	it is sufficient to take the empirical version of $Var[X_1] + Var[X_2]$.
	
\end{frame}

\end{document}
